\documentclass[a4paper,12pt]{report}
\usepackage[brazil, english]{babel}
\usepackage[utf8]{inputenc}
\usepackage[T1]{fontenc}
\usepackage{geometry}
\usepackage{setspace}
\usepackage{titlesec}
\usepackage{hyperref}
\usepackage{graphicx}
\usepackage{caption}
\usepackage{subcaption}
\usepackage{fancyhdr}
\setlength{\headheight}{15pt}
\addtolength{\topmargin}{-2.5pt}
\usepackage{xcolor}
\usepackage{amsmath, amssymb, bm}
\usepackage{mathtools}
\usepackage{cancel}
\usepackage{tikz}
\usepackage{newunicodechar}
\usepackage{ragged2e}
\usepackage{setspace}
\usepackage{tikz-3dplot} % Necessário para coordenadas 3D
\usetikzlibrary{intersections}
\usepackage{siunitx}
\usetikzlibrary{3d, arrows.meta}
\usepackage{booktabs}


\usepackage{color}
\definecolor{myblue}{rgb}{.8, .8, 1}

\usepackage{amsmath}
\usepackage{empheq}

\newlength\mytemplen
\newsavebox\mytempbox

\makeatletter
\newcommand\mybluebox{%
    \@ifnextchar[%]
       {\@mybluebox}%
       {\@mybluebox[0pt]}}

\def\@mybluebox[#1]{%
    \@ifnextchar[%]
       {\@@mybluebox[#1]}%
       {\@@mybluebox[#1][0pt]}}

\def\@@mybluebox[#1][#2]#3{
    \sbox\mytempbox{#3}%
    \mytemplen\ht\mytempbox
    \advance\mytemplen #1\relax
    \ht\mytempbox\mytemplen
    \mytemplen\dp\mytempbox
    \advance\mytemplen #2\relax
    \dp\mytempbox\mytemplen
    \colorbox{myblue}{\hspace{1em}\usebox{\mytempbox}\hspace{1em}}}
\makeatother

\usepackage[most]{tcolorbox}

\newtcbox{\mymath}[1][]{%
    nobeforeafter, math upper, tcbox raise base,
    enhanced, colframe=blue!30!black,
    colback=blue!30, boxrule=1pt,
    #1}

\tcbset{
    highlight math style={
        enhanced,
        colframe=red!60!black,
        colback=yellow!50,
        arc=4pt,
        boxrule=1pt,
        drop fuzzy shadow
    }
    }

\usepackage{physics}
\usepackage{pgfplots}
\pgfplotsset{compat=1.17}

\linespread{1.5}

\definecolor{ao(english)}{rgb}{0.0, 0.5, 0.0}
\definecolor{byzantium}{rgb}{0.44, 0.16, 0.39}
\newunicodechar{∘}{\circ}

\input{commands.tex}

\hypersetup{
    colorlinks=true,% make the links colored
    linkcolor=blue
}

\usepackage{setspace}
\addbibresource{bibliography.bib}

\newcommand{\printingbibliography}{%

    \pagestyle{myheadings}
    \markright{}
    \sloppy
    \printbibliography[heading=bibintoc, % add to table of contents
                   title=Refer\^encias % Chapter name
                  ]
    \fussy%
}
\PassOptionsToPackage{table}{xcolor}

\pagestyle{fancy}
\fancyhf{}
\renewcommand{\headrulewidth}{0pt}
\fancyhead[R]{\thepage}

\geometry{a4paper,top=30mm,bottom=20mm,left=30mm,right=20mm}

\titleformat*{\section}{\bfseries\large}
\titleformat*{\subsection}{\bfseries\normalsize}

\title{\LARGE \textbf{Preparation Plan -- Tech 3 (Master's Level)
\large Instituto Itaú de Ciência, Tecnologia e Inovação \\
\large Inova Talentos}}
\author{Andr\'e Vieira da Silva}
\date{\today}

\begin{document}

\noindent\textbf{Key Objective:} \\
Combine strong analytical skills in data with the practical ability to design and deploy AI agents that generate measurable business impact.

\vspace{0.5cm}

\section*{1. Understanding the ``Tech 3 -- Master'' Profile}

\begin{itemize}
\item[] Advanced analytical reasoning
\item[] Technical autonomy
\item[] Strategic business thinking
\item[] Ability to propose leverage points
\item[] Capability to design end-to-end solutions
\end{itemize}

\vspace{0.5cm}

\section*{2. Data Analysis \& Business Levers (Critical)}

\subsection*{Advanced SQL}

\begin{itemize}
\item[] CTEs
\item[] Window functions
\item[] Complex aggregations
\item[] Correlated subqueries
\item[] Query optimization
\item[] Dimensional modeling (Star Schema)
\end{itemize}

\subsection*{Business Metrics and Levers}

\begin{itemize}
\item[] Conversion funnels
\item[] Cohort analysis
\item[] Retention metrics
\item[] LTV (Lifetime Value)
\item[] CAC (Customer Acquisition Cost)
\item[] ROI (Return on Investment)
\item[] Price elasticity
\end{itemize}

\vspace{0.5cm}

\section*{3. AI Development -- Agents}

\subsection*{Core Knowledge}

\begin{itemize}
\item[] LLM fundamentals
\item[] RAG (Retrieval-Augmented Generation)
\item[] Tool integration
\item[] Orchestration
\item[] Memory strategies
\item[] Multi-agent systems
\item[] Observability and monitoring
\end{itemize}

\subsection*{Technical Foundations}

\begin{itemize}
\item[] Embeddings
\item[] Vectorization
\item[] Vector databases (FAISS, Pinecone)
\item[] Chunking strategies
\item[] LLM evaluation frameworks
\item[] Cost vs performance trade-offs
\end{itemize}

\vspace{0.5cm}

\section*{4. Generative AI Depth}

\begin{itemize}
\item[] Frontier models
\item[] Context windows
\item[] Hallucination causes
\item[] Alignment principles
\item[] Safety mechanisms
\item[] Temperature / top-p
\item[] Prompt chaining
\item[] Structured outputs
\item[] Ability to justify technical decisions
\end{itemize}

\vspace{0.5cm}

\section*{5. Large-Scale Data Systems (Desirable)}

\begin{itemize}
\item[] Partitioning strategies
\item[] Indexing techniques
\item[] Query execution plans
\item[] Distributed systems concepts
\item[] Spark fundamentals
\item[] ETL vs ELT
\item[] Data warehouse vs Data lake
\end{itemize}

\vspace{0.5cm}

\section*{6. 3-Week Intensive Study Plan}

\subsection*{Week 1 -- SQL \& Business Analytics}

\begin{itemize}
\item[] Advanced SQL practice
\item[] Data modeling
\item[] Business metrics deep dive
\item[] Churn case study
\end{itemize}

\subsection*{Week 2 -- Generative AI \& Agents}

\begin{itemize}
\item[] Build an agent with tool calling
\item[] Implement a RAG pipeline
\item[] Evaluate model outputs
\item[] Measure inference cost
\end{itemize}

\subsection*{Week 3 -- Architecture \& Strategy}

\begin{itemize}
\item[] Design an end-to-end solution
\item[] Address scalability concerns
\item[] Simulate a technical interview
\item[] Prepare a project pitch
\end{itemize}

\vspace{0.5cm}

\section*{7. Interview Preparation}

Be ready to answer:

\begin{itemize}
\item[] How would you measure the impact of an AI agent?
\item[] How do you ensure reliability?
\item[] How do you reduce hallucinations?
\item[] How do you scale an LLM-based system?
\item[] How do you justify cost vs benefit?
\end{itemize}

\end{document}